\documentclass[border=12pt]{standalone}
\usepackage[utf8]{inputenc}
\usepackage[T1]{fontenc}
\usepackage[american, siunitx]{circuitikz}

\sisetup{output-decimal-marker={,}}

\ctikzset{
    resistors/scale=0.85
}

\begin{document}
\begin{circuitikz}[font=\sffamily\small]

    % =========================================================================
    % 1. BLOCO ESQUERDO: ESP32 DEVKITC V4 (U1)
    % =========================================================================
    \draw[thick, fill=gray!5] (-5.2, -2.2) rectangle (0, 2.2);
    \node[align=center, font=\sffamily\bfseries] at (-3.1, 0.3) {U1\\ESP32};
    \node[font=\sffamily\scriptsize, text=gray!70!black] at (-3.1, -0.6) {(DevKitC V4)};

    % Alimentacao de U1 (+3,3V e GND)
    \draw (-2.6, 2.2) -- (-2.6, 3.0) node[vcc]{+3,3V};
    \node[below, font=\ttfamily\footnotesize] at (-2.6, 2.2) {3V3};

    \draw (-2.6, -2.2) -- (-2.6, -3.0) node[ground]{};
    \node[above, font=\ttfamily\footnotesize] at (-2.6, -2.2) {GND};

    % Pino de entrada analogica (VN / GPIO39)
    \node[left, font=\ttfamily\scriptsize] at (0, 0) {VN (GPIO39)};

    % =========================================================================
    % 2. BLOCO DIREITO: POTENCIOMETRO (POT1)
    % =========================================================================
    \draw (3.6, -2.4) node[ground]{}
        -- (3.6, -1.6)
        to[pR, n=POT1, l_=POT1] (3.6, 1.6)
        -- (3.6, 2.4) node[vcc]{+3,3V};

    % Conexao do cursor (SIG) para VN (GPIO39)
    \draw[->, >=stealth, thick] (POT1.wiper) -- (0, 0)
        node[midway, above, font=\ttfamily\scriptsize] {SIG};

\end{circuitikz}
\end{document}